\extrachap{Introduction}
\addcontentsline{toc}{chapter}{Introduction}

A quasi-plurisubharmonic potential describes a singularity through its growth near its poles. Its multiplier ideals describe the same potential through local integrability, while its Monge--Amp\`ere measure records a nonlinear concentration of curvature. Understanding the relation between these descriptions is a central problem of global pluripotential theory. This book studies that problem for Hermitian backgrounds, taking the theory of \emph{Singularities in global pluripotential theory} \cite{SGPT} as its starting point.

In a fixed K\"ahler class, all bounded potentials have the same Monge--Amp\`ere mass. The non-pluripolar mass of a singular potential can therefore be compared with a fixed total mass, and this comparison connects analytic approximation, envelopes and energy. A nonclosed Hermitian background changes this relation at its first step: the mass can vary even among smooth potentials. Ordinary singularity type, which identifies potentials with bounded difference, no longer determines a single mass. Envelopes and multiplier ideals remain available, but their connection to mass requires different arguments.

This loss of closedness is distinct from the geometry of the manifold. A non-K\"ahler manifold can carry closed positive currents, and a K\"ahler manifold can be equipped with a nonclosed Hermitian form. Let $X$ be a connected compact complex manifold of dimension $n\geq1$, and let $\theta$ be a smooth real $(1,1)$-form on $X$. A function $u$ is $\theta$-plurisubharmonic if it is locally the sum of a plurisubharmonic function and a smooth function and the current $\theta_u=\theta+\ddc u$ is positive. If $\theta$ is closed, this current is locally a complex Hessian. For a general background, a smooth local correction still makes $u$ plurisubharmonic, but leaves a smooth remainder in the current. Compactness, regularization and multiplier ideals thus survive locally; the new difficulties arise when these constructions are related globally.

\Needspace{8\baselineskip}
\section*{Closedness, mass and energy}

Let $\omega$ be a Hermitian form on $X$, that is, a smooth strictly positive $(1,1)$-form. If $\omega$ is K\"ahler, every smooth real potential $u$ with $\omega+\ddc u>0$ satisfies
\[
 \int_X(\omega+\ddc u)^n=\int_X\omega^n.
\]
The identity extends to bounded $\omega$-plurisubharmonic potentials. For singular potentials, the non-pluripolar product can have smaller total mass, and the difference measures the mass excluded at the poles.

For a general Hermitian form, integration by parts gives instead
\[
 \int_X(\omega+\ddc u)^n-\int_X\omega^n
 =\sum_{k=1}^{n-1}\binom nk
   \int_X u(\ddc u)^{k-1}\wedge\ddc(\omega^{n-k})
\]
for every smooth $u$. The derivatives of the background already affect the mass before singularities appear. \Cref{chap:mass} proves that mass invariance is equivalent to $\ddc(\omega^j)=0$ for $1\leq j\leq n-1$. The usual global comparison inequality would force this invariance, so a comparison argument in an arbitrary Hermitian background must account for these derivatives.

Even an upper or a positive lower mass bound is an additional condition. The scalar Hopf manifolds of dimension at least three have unbounded masses of smooth potentials with positive perturbed form, while the standard Hopf threefold has a positive infimum of those masses. A different compact threefold carries a balanced metric with bounded masses whose infimum is zero (\cref{chap:hopf,sec:vanishing-lower-mass}). The two bounds therefore cannot replace one another. The general constructions of the book are made without either bound; \cref{part:mass-bounds} develops the further theory that uses them.

Mass conservation itself is weaker than the closedness needed for the K\"ahler energy. At a smooth potential $u$ with $\omega+\ddc u>0$, consider the one-form
\[
 \alpha_u(v)=\int_X v(\omega+\ddc u)^n
\]
on the space of such potentials, where $v$ is a smooth real tangent direction. For a K\"ahler form, it is the differential of the Monge--Amp\`ere energy. In dimension at least two, a primitive even near the zero potential forces $\mathrm d\omega=0$ (\cref{thm:raw-primitive}). On a non-K\"ahler surface, a Gauduchon metric conserves mass but has no such primitive. Perron interpolation and Legendre duality can still be constructed by envelopes; an interpretation through a metric energy needs a separate justification.

A linear conservation law supplies another form of global control. A Gauduchon form $\gamma$ satisfies $\ddc(\gamma^{n-1})=0$, and hence, for every $\theta$-plurisubharmonic potential,
\[
 \int_X(\theta+\ddc u)\wedge\gamma^{n-1}
 =\int_X\theta\wedge\gamma^{n-1}.
\]
This trace measure determines $u$ up to an additive constant. In a fixed background, for a smooth positive volume form $\mathrm dV$ and the normalization $\int_Xu\,\mathrm dV=0$, weak convergence of trace measures is equivalent to $L^1$ convergence of potentials (\cref{chap:compactness}). The trace stays finite even when the nonlinear Monge--Amp\`ere mass is infinite. Extended non-pluripolar products and finite damped measures make it possible to study such potentials, while passage from damped to undamped measures requires control near the poles.

\section*{Singularity types and envelopes}

Multiplier ideals record the holomorphic functions locally square integrable against $\mathrm e^{-ku}$, for every $k>0$. This gives a second equivalence relation on potentials, coarser than ordinary singularity type. To compare the two, fix $u\in\PSH(X,\theta)$ and take envelopes in the same background. The ordinary envelope $P_\theta[u]$ is the upper semicontinuous regularization of the supremum of the $\theta$-plurisubharmonic $h\leq0$ with $h\leq u+C$ for some constant $C$. The multiplier-ideal envelope $P_\theta[u]_{\mathcal I}$ uses instead the $h\leq0$ with the same multiplier ideals as $u$ at every positive weight.

In a closed big class on a compact K\"ahler manifold, for potentials of positive mass, the equality of these two envelopes characterizes $\mathcal I$-good singularities and is equivalent to a mass identity and approximation by analytic singularities. In the Hermitian setting, analytic approximation still relates the envelopes even when a mass identity is unavailable. Correcting the curvature errors keeps the dyadic approximants in a positive Hermitian background, and their envelopes decrease to the multiplier-ideal envelope. In such a positive background, comparing the two envelopes defines a uniform topology on ordinary classes in which their equality characterizes the closure of the analytic classes, including classes with infinite-mass representatives (\cref{thm:ordinary-topology-analytic-closure}). This topology can identify distinct ordinary classes.

The obstacle in an envelope is part of this construction. A smooth change $\theta\mapsto\theta+\ddc\chi$ transforms potentials by $u\mapsto u-\chi$ and transforms their obstacles as well. Keeping only the obstacle zero need not give an invariant topology; allowing all smooth obstacles does. Whether equality of the envelopes at zero implies equality at every smooth obstacle is therefore a substantive question. It holds in a smooth big background with bounded mass when the multiplier-ideal envelope has positive lower volume, as well as on curves and on certain product loci without those volume conditions (\cref{prop:conditional-i-characterization,cor:curve-arbitrary-background-criterion,cor:general-auxiliary-all-obstacle-criterion}). The general case at zero volume remains open.

The mass can be attached to an ordinary type by taking its infimum and supremum over all representatives. These lower and upper singularity volumes need not agree. An ordinary model is a fixed point $P_\theta[u]=u$. In a smooth big background with bounded mass, volume comparison gives a completely metrizable topology on the models of positive lower volume. At zero volume, the same comparison can fail to distinguish models, so a topology extending this theory must keep information that mass has lost.

Curves show both the possibility and the limitation of a different construction. A curve model is determined by its polar curvature, the part of $\theta+\ddc u$ carried by the pole set. Total variation of that curvature gives a complete metric including zero-volume models, with the ideal models as the closure of the analytic models. It also gives complete loci on higher-dimensional products without an ambient mass bound (\cref{prop:general-auxiliary-product-metric}). In higher dimension, however, residual top-degree mass may lie on a diffuse pole set although all polar curvature currents vanish after every modification. Even the residual measure and both volumes can coincide for distinct models, while their comparison distance is positive (\cref{prop:equal-residual-anisotropic-models}). These phenomena already occur in closed K\"ahler backgrounds and delimit what a classification by measures can recover.

Restriction raises a related question. A potential may be identically $-\infty$ on a submanifold even though its generic Lelong number there is zero. In this case, the restrictions of its ambient multiplier ideals define finite divisorial data on every modification. Trace envelopes turn this data into constraints on potentials of the submanifold. Exact realization is possible on curves and in sufficiently large backgrounds on K\"ahler targets, and projective pullback data also admit exact realizations on Hopf manifolds of unbounded mass (\cref{thm:hopf-projective-exact-trace}). Thus failure of a mass bound does not itself prevent realization. Exact divisorial data nevertheless need not determine the Monge--Amp\`ere mass, or a single ordinary type.

\section*{Equations with prescribed data}

The varying mass also changes the normalization of an equation. Given a smooth real function $F$ on $X$, one seeks a smooth pair $(u,c)$ satisfying
\[
 (\omega+\ddc u)^n=c\,\mathrm e^F\,\mathrm dV,
 \qquad \omega+\ddc u>0,\qquad \sup_Xu=0.
\]
In a K\"ahler background, integration determines $c>0$ before solving for $u$. In a Hermitian background the left-hand mass depends on the unknown potential, so the constant must be found together with it. Such a smooth pair exists and is unique. The linearization gives an adjoint Gauduchon probability measure describing the variation of $\log c$ with the density, and leads to a variational volume defined by Gauduchon test forms. This volume remains finite even on manifolds with unbounded upper Monge--Amp\`ere mass.

Positive variational volume, however, need not imply that a background is big, meaning that it admits a potential $\rho$ with $\theta+\ddc\rho\geq\delta\omega$ for some $\delta>0$. A nef nonclosed background on a projective manifold has positive variational volume and fails this condition (\cref{thm:radial-positive-volume-nonbig}). Up to a smooth change of representative, the same example occurs at a finite maximal Chern--Ricci time (\cref{thm:actual-flow-positive-volume-nonbig}). An argument for geometric convergence that needs bigness therefore cannot substitute positive variational volume for it.

For a nonnegative, nonzero $L^p$ density, $p>1$, continuous solutions still exist and the positive constant is still unique, but uniqueness of the potential is a separate matter. A smooth density with zeros may have distinct normalized smooth solutions. For a continuous density whose positive set is connected and dense, bounded solutions are unique up to constants (\cref{thm:connected-positive-uniqueness}). Selection by approximation must also be distinguished from uniqueness. Moreover, when singularities are prescribed, distinct divisorial types can have the same normalization functional and supporting Gauduchon probabilities (\cref{chap:normalization}); the scalar response does not recover the full type.

Equations in a prescribed ordinary type use conditions on the reference potential and relative envelope estimates to keep the difference bounded. Equations with a general finite non-pluripolar measure need other existence hypotheses. \Cref{part:mass-bounds} treats semipositive big backgrounds of positive lower mass, with a model reference admitting a competitor in a smaller background, and arbitrary smooth big backgrounds of bounded mass, with a model reference of positive lower volume. The solution obtained belongs to a relative envelope class; bounded difference from the reference requires a further argument. Keeping these two conclusions separate is essential when using the equation to study singularities.

\section*{Geometric constructions}

Logarithmic poles give concrete tests of the theory. On a resolution, removing a prescribed divisorial pole can leave a strictly positive residual form, where a bounded equation determines the remaining potential. A divisor forced by the background behaves differently: it belongs to the type of minimal singularities, and removing it can force the current to descend. The constructions in \cref{sec:prescribed,sec:blowups} distinguish a freely prescribed pole from such a compulsory one.

The regularity of a potential also affects the geometry of its current. When that current is smooth and strictly positive off an analytic set, a H\"older potential places the analytic set at finite distance, making the regular metric incomplete if the set is nonempty (\cref{thm:holder-current-distance}). Boundedness alone is insufficient: radial examples in closed K\"ahler classes have continuous minimal residuals and complete regular ends. The integral controlling the potential differs from the one controlling path length. This distinction comes from regularity, even before closedness is lost.

Holomorphic sections link singularities to geometry through a closed line-bundle curvature, so counting sections cannot measure an arbitrary Hermitian Monge--Amp\`ere mass. Gauduchon degrees control section growth and flag valuations, while fibrations can concentrate equilibrium measures on a lower-dimensional base. Partial Bergman kernels connect these section spaces to envelopes of singular weights. Closed Bott--Chern classes and locally conformally K\"ahler structures give further settings where the differential constraints can be used more precisely.

\section*{Reading the book}

We assume familiarity with complex manifolds, differential forms and currents. \Cref{chap:framework} develops the plurisubharmonic tools used throughout. On a first reading, the maximum principle, normalized compactness and exponential integrability are the starting points; the longer approximation and descent arguments can be consulted when needed. The mass and energy calculations explain the obstructions above, and Gauduchon traces, non-pluripolar products and the Hopf construction complete the foundations.

The chapters on envelopes develop singularity comparison, analytic approximation and interpolation, while the equation chapters explain which estimates keep a prescribed type. Restriction and holomorphic sections then connect the analytic constructions to geometry. \Cref{part:mass-bounds} develops volume comparison and equations with general measures under explicit mass hypotheses. The appendix treats potentials and equations on compact unibranch spaces, including projective cones. \Cref{chap:research-program} collects open questions, and \cref{chap:history} explains the sources and precise dependencies; monograph references also appear where they are used.
